% Transcription — fonds Alexandre Grothendieck, shelfmark 156-3, batch 2 (pages 21-40).
% Established from the facsimile archives/batches/156-3/batch-02.pdf (a mirror of
% https://grothendieck.umontpellier.fr/156-3.pdf; no cover sheet in the batch, so
% PDF page k = folder page 20+k), per the skill transcribe-grothendieck. The
% folder is forty pages; this is the second and last batch. Batch 1 (pages 1-20)
% is transcribed separately; its last pages (read here only for context) define
% « parties modérées », the relative position « modérée » (pos. rel. mod.) of
% two parts of a tronçon, and C_mod(Y,X), and end on the least closed modérée
% part Z containing C_mod(Y,X), which page 21 takes up.
%
% All twenty leaves are mathematics and all are transcribed; nothing is skipped
% and nothing is another's text. His own pagination 21-40 runs across the top of
% the leaves and coincides with the archivists'. One date in his hand:
% « (9 juin) » on page 21, opening the section « ① Topologie sur L ». The
% inventory's bracketed « [Chapitre] » is the archivists'.
%
% What the batch is about. Page 21, first third: the end of batch 1's argument,
% C_mod(Y,X) = C_mod(F,Z) with F the relative frontier of Y in X. Then (9 June)
% a topology on a « tronçon ordonné » L, with basis of opens ]a,b[, L_{<a},
% L_{>b}: points need not be closed; page 22 notes the space is T_0, asks
% whether it is sober, and observes that only for L totally ordered are points
% closed; a second topology is proposed with basic closed sets [a,b], L_{>=a},
% L_{<=a}. Page 23: in it ]a,b[ need no longer be open, the coarser common
% topology, and the questions whether the closure of ]a,b[ is [a,b] (false
% already for the set of lieux of a faille topologique, margin). Pages 24-26:
% he settles on the first topology, calls [a,b], L_{<=a}, L_{>=b}, L « closed
% sub-tronçons », defines the absolute boundary ∂T (extreme points), proves
% ∂T = (Fr T ∩ T) ⊔ (T ∩ ∂L), then tronçons T' = T \ α (α ⊂ ∂T), open and
% semi-open ones, « pseudo-tronçons » (tronçon or point), and two
% pseudo-tronçons « en position relative modérée » when ∂T ∪ ∂T' is totally
% ordered. Pages 27-28: T ∩ T' is again one; the order T ≺ T' (every t < t'),
% « raccordement », the strict order T ≪ T', and the interstitial
% pseudo-tronçon between T ≪ T' (Prop., Cor. 1-2: left and right interstitials).
% Pages 29-31: Tr(L) ordered by ≪; for a finite chain Φ, X_Φ = ⋃ T_i; Φ ↦ X_Φ
% injective on flags; X_Φ are the « modérées » parts, with Fr(X_Φ), components
% π_0(X), the complement X' = C_mod(X) described by the interstitials, and
% |card π_0(X) − card π_0(X')| ≤ 1 with the three cases. Pages 31-33: the
% structure on M = Mod(L): order, the relation {X,Y} modérée, closed modérées
% M_f; axioms 1)-8) (least/greatest element, inf/sup of a modérée pair, the
% complement X' as an involutive anti-isomorphism, closure X̄, frontier, « rare »
% elements), and the Boolean-ring notation X+Y, X·Y. Pages 34-36: relativisation
% to M_A = M_{≤A}; he checks 1)-7) for M_A one by one. Pages 36-38: connexity
% (closed and open ⇒ ∅ or 1); for A = X_Φ the connected components are the
% pseudo-tronçons T_i. Page 39: the Boolean algebra generated by a finite family
% of mutually modérées parts is determined by the chain S of their frontier
% points. Page 40: « Caractérisation d'un intervalle ordonné, en tant que
% réseau », by S ⊂ 𝔓_2(L) and the decomposition of ∁_mod ε; he stops, thinking
% it must have been done long ago (« Hilbert ? »).
%
% Reading hazards fixed once here. L is his cursive L for the tronçon (plain $L$
% here); the X he often writes with an extra oblique stroke, transcribed as a
% plain $X$; the family of modérées parts is 𝓜 (\mathcal{M}); on pages 29-30
% he writes U (\mathcal{U}) for X_Φ interchangeably with X. The dotted Ṫ, Ẋ are
% frontiers (Fr). « ps.-tr. » = pseudo-tronçon, « pos. rel. mod. » = position
% relative modérée, « OPS » = on peut supposer. His underlining is set in bold.
% The hand is fast drafting throughout; page 21's second half, the margins of
% pages 26, 30 and 32, and the NB of pages 36 and 39 are overwritten and only
% partly read.
%
% Readings flagged and not settled: « sobre » (p. 22); the second list of basic
% opens on page 23 (∁L_{≥a}, ∁L_{≤a}, [a,b] as written); « ω(L) » for the
% greatest element (p. 24); the formula for ∂T' on page 26; « l'ens. des » in
% Cor. 1 (p. 28); what stands under Drap on page 29 (« y inclus … »); the
% interlinear word under « Ti ∩ A = Ti » (p. 38). Variances on the page and not
% repaired: « [b,a[ » for [a,b[ (p. 26); « Fr X ∪ Fr X » (p. 33); « 1_{M_A} = X »
% for A and « Y ∈ M_Y » for M_A (p. 34); « ]t_1, t_3[ » in the list of page 39.
% Page 25's « ∂L = ∂L (sic) » is his own « sic ».
%
% Pass: Opus 5.5 (claude-opus-5-5), 2026-09-24 — first pass, unchecked against the pages by a human.
\documentclass[11pt,a4paper]{article}
% Le préambule commun à toutes les transcriptions du fonds Grothendieck.
%
% Il tient en une idée : **l'incertitude se compose, elle ne se résout pas.**
% Une transcription qui lisse un mot douteux a détruit la seule chose pour
% laquelle on la faisait — un lecteur doit pouvoir savoir, à chaque endroit,
% ce qui était sur le feuillet et ce qui a été ajouté. Les six macros qui
% suivent sont donc l'appareil critique, et non de la décoration.
%
% Les mêmes commandes sont reconnues par `scripts/render.mjs`, qui produit la
% vue de lecture du site. Une macro ajoutée ici doit l'être aussi là-bas,
% faute de quoi le rendu échouera bruyamment — ce qui est voulu : un rendu
% silencieusement faux est pire qu'un rendu absent.

\ProvidesPackage{grothendieck}[2026/08/08 Transcriptions du fonds Grothendieck]

\usepackage{amsmath,amssymb,amsthm}
\usepackage{mathrsfs}
\usepackage{stmaryrd}
% Les diagrammes commutatifs sont partout dans ce fonds : c'est la langue
% même de la géométrie algébrique de Grothendieck, et une transcription qui
% les rendrait en prose ne transcrirait rien.
\usepackage{tikz-cd}
% Français par défaut — c'est la langue des feuillets — mais l'anglais est
% chargé aussi : la traduction et le résumé sont des documents anglais, et la
% typographie française (espaces avant les deux-points) y serait une faute.
% Ces documents basculent par \selectlanguage{english} en tête de corps.
\usepackage[english,french]{babel}
\usepackage{fontspec}
\usepackage{geometry}
\geometry{a4paper,margin=25mm,marginparwidth=28mm}
\usepackage{marginnote}
\usepackage{xcolor}
% ulem, not soul. `soul`'s \st and \ul are built on a character-by-character
% scan that XeLaTeX's font machinery breaks: the document fails with an
% undefined control sequence rather than degrading. ulem does the same two jobs
% and is Unicode-engine safe. `normalem` stops it hijacking \emph, which the
% transcriptions use for Grothendieck's own emphasis.
\usepackage[normalem]{ulem}
\usepackage{hyperref}
\hypersetup{colorlinks=true,linkcolor=black,urlcolor=black,pdfborder={0 0 0}}

% --- Métadonnées du lot -----------------------------------------------------
% Reprises telles quelles par le script de rendu, qui les affiche en tête de
% la vue de lecture. Elles fixent ce que le fichier prétend couvrir : sans
% elles, une transcription flotte sans point d'attache dans un fonds de seize
% mille feuillets.

\newcommand{\folder}[1]{\gdef\grothFolder{#1}}
\newcommand{\batch}[1]{\gdef\grothBatch{#1}}
\newcommand{\leaves}[2]{\gdef\grothFirst{#1}\gdef\grothLast{#2}}
\newcommand{\foldertitle}[1]{\gdef\grothFoldertitle{#1}}
\gdef\grothFolder{?}\gdef\grothBatch{?}\gdef\grothFirst{?}\gdef\grothLast{?}\gdef\grothFoldertitle{}

% --- Le résumé --------------------------------------------------------------
%
% Il ouvre la lecture modernisée, sous le titre « Résumé », et il est composé
% en petit. Ce n'est pas un ornement : c'est le seul passage du document qui
% ne soit pas des mathématiques, et un lecteur doit pouvoir le reconnaître
% avant de l'avoir lu — puis le sauter s'il connaît déjà le sujet, ou s'y
% arrêter s'il ne le connaît pas. Le filet qui le suit marque la reprise du
% corps.

\newenvironment{resume}
  {\small\setlength{\parskip}{0.45em}}
  {\par\medskip\hrule\medskip}

% --- Mots-clés ---------------------------------------------------------------
%
% En anglais, en clôture du résumé de la lecture modernisée : le vocabulaire
% moderne sous lequel chercher ce que les feuillets construisent. Le manifeste
% du site (`npm run manifest`) les extrait pour étiqueter la cote entière —
% cette ligne est la source unique des tags, il n'y a pas de fichier à côté.
\newcommand{\keywords}[1]{\par\smallskip\noindent
  {\footnotesize\textsc{Keywords} --- \textit{#1}}\par}

% --- L'appareil critique ----------------------------------------------------

% Le numéro de feuillet, en marge. C'est l'ancre qui permet de régler un
% désaccord sur une formule en regardant *un* feuillet plutôt qu'une chemise
% de six cents. Le site s'en sert aussi pour tourner les pages du fac-similé
% quand on fait défiler la transcription.
\newcommand{\leaf}[1]{%
  \marginnote{\footnotesize\color{gray!70}\textsf{#1}}%
  \label{leaf:#1}\ignorespaces}

% Illisible : on ne devine pas. Un mot inventé qui se lit comme les autres est
% le pire résultat possible.
\newcommand{\ill}{\textcolor{red!60!black}{[\,\ldots\,]}}

% Lecture douteuse : le mot est donné, et signalé comme incertain.
\newcommand{\uncertain}[1]{\uline{#1}}

% Ajout de l'éditeur — une lettre manquante, un mot sous-entendu.
\newcommand{\add}[1]{\textcolor{blue!50!black}{[#1]}}

% Biffé par Grothendieck lui-même. Ce qu'il a rayé fait partie du document :
% une rature dit souvent où la pensée a changé de direction.
\newcommand{\struck}[1]{\sout{#1}}

% Note du transcripteur, sur l'état matériel du feuillet.
\newcommand{\note}[1]{\par\smallskip{\small\color{gray!60!black}\textit{[#1]}}\par\smallskip}

% Note marginale de Grothendieck, distincte du corps du texte.
\newcommand{\marginal}[1]{\par\smallskip\noindent\hspace{1em}%
  {\small\color{gray!60!black}#1}\par\smallskip}

% --- Titre ------------------------------------------------------------------

\AtBeginDocument{%
  \begin{center}
    {\large\bfseries Fonds Alexandre Grothendieck --- cote n\textsuperscript{o}~\grothFolder}\\[2pt]
    {\itshape\grothFoldertitle}\\[4pt]
    {\small feuillets \grothFirst--\grothLast\ (lot \grothBatch)}\\[2pt]
    {\footnotesize\color{gray!60!black}Transcription établie d'après le fac-similé numérisé
     par l'Université de Montpellier.\\
     Les lectures incertaines sont soulignées, les illisibles notées [\,\ldots\,],
     les ajouts entre crochets.}
  \end{center}
  \vspace{1em}}

\endinput


\folder{156-3}
\batch{2}
\pages{21}{40}
\dating{1986}
\watermark{Édition de démonstration}
\foldertitle{[Chapitre] III. Réseaux via découpages : notes manuscrites (08/06/1986).}

\begin{document}
\page{21}
\note{la page continue l'énoncé de la fin de la page 20 : le plus petit ensemble fermé modéré $Z$ contenant $C_{\mathrm{mod}}(Y,X)$}

Si on pose $F$ = frontière relative de $Y$ dans $X$, \struck{(F} Alors $F$ est fermé, \add{modéré,} en pos. mod. par rapport à $Y$, \struck{donc} $F \subset Z$, \struck{\ill{}} \ill{} $F$ \uncertain{ne} \struck{\ill{}} contient pas de pts isolés de $Z$.
\note{le $F$ de la première ligne et celui de la formule ci-dessous sont tracés d'une barre verticale en surcharge ; le mot « modéré » est inséré au-dessus de la ligne}
\struck{Comp}
\[
  C_{\mathrm{mod}}(Y,X) = C_{\mathrm{mod}}(F, Z)
\]

\subsection{(9 juin)}

Soit $L$ un tronçon ordonné.

\struck{Un sous-ensemble \ill{} dans sous-tronçon fermé de $L$ est un sous-ens. de la forme $[a,$}
\note{cette définition inachevée est encadrée et barrée de plusieurs traits obliques}

\textbf{(1) Topologie sur $L$} : base d'ouverts les $]a,b[$ (pour $a < b$), $L_{<a}$, $L_{>b}$

[NB si $L$ n'a pas de plus petit élément (resp. de plus grand élt), on peut omettre les $L_{<a}$ (resp. les $L_{>b}$).
\struck{Une} caractérisation \ill{} \struck{Pour que \ill{} soit} \ill{} (s'il y en a) est fermé, mais \struck{fermé, il faut \ill{} que} \uncertain{OK} \ill{}.
\ill{} qu'on \uncertain{généralise} ?, \struck{\ill{}} \struck{\ill{}} pour $a \in L$, $\{a\}$ n'est pas \add{néc.} fermé. Si $L$ est l'un des lieux d'une \struck{top} \ill{} \ill{} \uncertain{Alors} $b \in \overline{\{a\}}$ i.e. $\overline{\{b\}} \subset \overline{\{a\}}$, \uncertain{équivalant} à $b \supset a$ \ill{}]
\note{ces lignes sont très surchargées ; seules les formules sont sûres}

\page{22}
Il en résulte que $\overline{\{a\}} = \overline{\{b\}} \Leftrightarrow a = b$. Cela exprime que la top. $L$ satisfait l'axiome $T_0$, il est « \uncertain{prisotope} » i.e. s'injecte dans le \uncertain{prisotopifié}) — je ne sais si c'est le cas pour tout tronçon, ni si un tronçon, avec sa top. canonique, est \uncertain{sobre} (donc, de plus, toute partie fermée irréductible aurait un pt générique…) Ainsi, les « segments » $[a,b]$ ne sont \uncertain{pas néc.} fermés pour la top. canonique. Ce n'est que lorsque $L$ est \textbf{tot} ordonné qu'on peut conclure que les pts sont fermés (réciproque ?) et en fait, $L$ est \textbf{séparé}, et de plus les $[a,b]$, $L_{\geq a}$, $L_{\leq a}$ sont fermés.

On peut définir une autre topologie, ayant comme base fondamentale de fermés les $[a,b]$, $L_{\geq a}$, $L_{\leq a}$ (\uncertain{étant} entendu qu'on prend les intersections finies des réunions finies, pour avoir tous les fermés)

\page{23}
Mais maintenant, on est ennuyé, car $]a,b[$ n'a plus aucune raison d'être ouvert !

\struck{Ici} On peut bien sûr prendre la topologie \uncertain{la moins fine} des deux, ayant comme base d'ouverts
\[
  \begin{array}{ll}
    L_{<a},\ L_{>a},\ ]a,b[ & (a < b) \\
    \complement L_{\geq a},\ \complement L_{\leq a},\ [a,b] & (a \leq b)
  \end{array}
\]
\note{ainsi sur la page ; on attendrait $\complement [a,b]$ dans la seconde ligne}

\struck{Dire que $]a,b[$ est \ill{}} \struck{Est-il vrai, si $a \leq b$, et} \struck{Mais} \struck{$b$ n'est pas plus grand élément, que} \struck{Avec \ill{} $b \in [a,b$}
\[
  \overline{]a,b[} = [a,b] \ ?
\]
\[
  \overline{L_{<a}} = L_{\leq a} \ ? \quad \text{(si $a$ n'est pas plus petit élément ($a \neq 0_L$))}
\]
\[
  \overline{L_{>a}} = L_{\geq a} \ ? \quad \text{(si $a$ n'est pas plus grand élément ?)}
\]

\struck{Pour} Le cas \uncertain{crucial}, dans le premier cas, est de prouver (\uncertain{i.e.}) $a \in \overline{]a,b[}$, et pour cela, il suffit de voir \struck{que} qu'un ouvert « élémentaire » qui contient $a$, contient un $]a, \varepsilon[$, avec $a < \varepsilon < b$. Mais c'est faux déjà pour un ens. de lieux associé à une faille topologique, \uncertain{cf.} [\uncertain{pratiquement}] él. $\complement[x,y]$

\marginal{\ill{} \ill{} que $a \in \complement\,[x,y]$ mais $]a,b[ \subset [x,y]$}
\note{la marge porte un croquis : une bande hachurée sur laquelle sont marqués $x$, $a$, $b$, $y$, la partie entre $a$ et $b$ plus fortement hachurée}

\page{24}
Finalement, on travaillera avec la top. définie d'abord (base \uncertain{formée} \ill{} les « intervalles »), et on appellera néanmoins \uncertain{par abus de langage}, les $[a,b]$, $L_{\leq a}$, $L_{\geq b}$ (ainsi que $L$) \textbf{sous-tronçons} ordonnés (ou biordonnés) « fermés », [NB on les suppose non réduits à un pt, i.e. $a < b$ resp. $a \neq 0(L)$ resp. $b \neq \omega(L)$]
\note{la lettre lue $\omega$ est douteuse ; elle désigne le plus grand élément de $L$}

Un \textbf{sous-tronçon \add{fermé} \struck{large}} est une \struck{sous-tronçon} partie qui est soit un ss-tronçon, soit réduite à un pt. Pour une telle partie $T$, on \struck{pose}
$\partial T$ = \struck{$\{ t \in T \mid t \in \ldots \}$ \ill{}} bord \uncertain{absolu} de $T$ = ensemble \struck{canonique} fini de ses extrémités \struck{$= \overline{T} \cap \complement T$}
(NB $\mathrm{card}\, \partial T \in \{0,1,2\}$)

On voit que pour un tronçon fermé (NB non réduit à un pt) on a
\[
  \partial [a,b] = \{a,b\}, \quad \partial L_{\leq a} = \{a\}, \quad \partial L_{\geq a} = \{a\}, \quad \partial L = \emptyset
\]
\note{tout ce passage, depuis « on pose », est encadré et traversé de traits obliques ; il est repris à la page suivante}

\page{25}
pour
\[
  \partial T = \text{bord absolu de } T = \Bigl\{ t \in T \;\Big|\; t \text{ est un plus petit ou plus grand élément de } T \Bigr\}
\]

\struck{Si $T$ est réduit à un pt,} On a donc
\[
  \begin{array}{ll}
    \partial \{a\} = \{a\} & \\
    \partial [a,b] = \{a,b\} & \text{si } a \leq b \\
    \partial L_{\geq a} = \partial L_{\leq a} = \{a\} & \forall a \in L \\
    \partial L = \partial L \ \text{(sic)} &
  \end{array}
\]
\note{devant la dernière ligne, un mot biffé et illisible ; le « (sic) » est de sa main}
\marginal{(\ill{} \ill{} \ill{} $= \partial L$}

Notons que
\[
  \dot{T} \cap T \subset \partial T \subset (\dot{T} \cap T) \cup \partial L
\]
i.e. tout $t \in T$ qui est pt frontière, est dans le bord \add{de $T$}, et l'inverse est vrai si $t \notin \partial L$. Plus précisément
\struck{si $\mathrm{card}\, T = 1$, $\partial T = \dot{T} = T$ ; sinon $\dot{T} \cap T = \partial T \smallsetminus (\partial T \cap \partial L)$, $\partial T = (\dot{T} \cap T) \cup T \cap \partial L$}
\note{ce premier énoncé, encadré, est barré de traits obliques et remplacé par la formule encadrée qui suit}
\[
  \boxed{\partial T = (\dot{T} \cap T) \sqcup (T \cap \partial L)}
\]

\textbf{Tronçons} : C'est \uncertain{une} partie \struck{\ill{} réduite \ill{}} de $T$ \ill{} qui forme $T' = T \smallsetminus \alpha$, avec $\alpha \subset \partial T$. \struck{\ill{}} \struck{Tronçons} \ill{} un ens. de la forme
\[
  \left\{
  \begin{array}{ll}
    ]a,b],\ [a,b[,\ ]a,b[ & (a < b) \\
    L_{<a} & (a \neq 0_L) \\
    L_{>a} & (a \neq \omega_L) \\
    L &
  \end{array}
  \right.
\]
\note{après $L_{<a}$, une surcharge biffée illisible}
\marginal{NB Cette écriture est \textbf{unique} (on pose $T = \overline{T'}$ : c'est une adhérence, le plus petit fermé contenant $T'$)}

\page{26}
Les tronçons $]a,b[$, $L_{<a}$, $L_{>a}$, $L$ sont dits \textbf{ouverts} \struck{(NB ils sont} (ce sont bien ceux qui sont ouverts pour la top. can.), \struck{(ce sont bien ceux} et $]a,b]$ \struck{qui sont ouverts} et $[b,a[$ sont dits \textbf{semi-ouverts}.
\note{« $[b,a[$ » sur la page ; on attend $[a,b[$}

\textbf{Attention} \struck{le tronçon} Si $a$ plus petit élément, alors $[a,b[$, \struck{\ill{}} semi-ouvert pour la terminologie précédente, est aussi ouvert pour la top. canonique, et il est égal à $L_{<b}$. Kif kif dualement. $L$ est à la fois ouvert et fermé.

Un \textbf{tronçon généralisé} \add{ou ps.-tronçon} est une partie de $L$ qui est soit réduite à un pt, soit un tronçon. Pour un tronçon de la forme
\[
  T' = T \smallsetminus \alpha \qquad (T \text{ tronçon fermé, } \alpha \subset \partial T)
\]
\struck{\ill{}} on pose
\[
  \partial T' = \partial T
\]
\[
  \partial T' = \dot{T}' \sqcup \partial L \ \ldots
\]
\note{lecture douteuse de cette seconde formule}
\marginal{Attention, cette fois ce n'est pas le $\partial$ absolu de $T'$ pour la structure d'ordre induite.}
\marginal{NB $\partial T \neq \emptyset$ sauf si \struck{\ill{}} $T' = L$, $L$ non \ill{} \ill{} \ill{} les $T'$ \ill{} (cas $T'$ ouvert) (cas $T'$ semi-ouvert) (cas $T'$ \ill{}) \ill{}}

\struck{Soient deux} \textbf{Deux} ps.-tronçons $T$, $T'$ sont dits \textbf{en pos. relative modérée} (abrég. $(T,T')$ modérée ou $\{T,T'\}$ \uncertain{modérée}) si $\partial T$ et $\partial T'$ le sont, i.e. $\partial T \cup \partial T'$ est tot. ordonné.
\marginal{\ill{} bornes de $T$ \ill{} bornes \ill{} $\partial T \cup \partial T'$}

\page{27}
\textbf{Prop} Si $T$, $T'$ sont des ps.-tr. en position rel. modérée, alors $T \cap T'$ est \add{vide ou est} un ps.-tr., et on a
\[
  \partial (T \cap T') \subset \partial T \cup \partial T'
\]

Immédiat. On va introduire une relation (pas d'inclusion !) d'ordre entre ps.-tronçons. On dit que
\[
  T \prec T'
\]
si \struck{a) $T$ est en pos. rel. modérée \ill{} $T'$}
\[
  \forall t \in T,\ t' \in T' \Longrightarrow t < t'
\]
\marginal{\struck{Si c'est} Si il faut que $T$ et $T'$ en pos. rel. mod. et $T \cap T' = \emptyset$}
C'est évidemment transitif.

Soit $T \prec T'$. Comme $T$ est majoré, il s'ensuit que \struck{$\partial T \neq \emptyset$} $T$ a un plus grand élément $a$, \struck{tel que} $T \subset L_{\leq a}$. \struck{Comme} De même, $\partial T'$ a un plus petit élément $b$, donc que $T' \subset L_{\geq b}$, \struck{\ill{}} et on voit de suite que $a \leq b$.

\textbf{Position relative \add{de raccordement}} \struck{impropre} : $a = b$, et $a \in T$ ou $b \in T'$ (à l'exclusion l'un de l'autre, puisque $T \cap T' = \emptyset$)

\struck{Dans la suite, on l'excluant de} On renforcera la relation $\prec$ en $\ll$ :
\[
  T \ll T' \iff T \prec T', \text{ et position relative propre}
\]
\[
  \iff \exists x \in L \text{ tel que } T \subset L_{\leq x},\ T' \subset L_{>x}
\]
\note{l'indice de $L_{\leq x}$ est surchargé ; on attend $L_{<x}$, comme à la page suivante}

\page{28}
\subsection{psTronçon interstitiel}

\textbf{Prop} Soient $T$, $T'$ \struck{ps.-tr. tels que $T \prec T'$, $a = \mathrm{Sup}\, T$, $b = \mathrm{Inf}\, T'$,} Conditions équivalentes (si $a = \mathrm{Sup}\, T$, $b = \mathrm{Inf}\, T'$) \struck{si $a \in$} :

a) $T \prec T'$ \add{et $T$, $T'$} ne se raccordent pas, i.e.
\[
  (a \in T \text{ ou } b \in T') \Longrightarrow a \neq b
\]

\struck{b) $\exists x \in L$} b) $\exists x \in L$ tel que $T \subset L_{<x}$, $T' \subset L_{>x}$
\marginal{NB b) implique que $T \prec T'$}

c) $\exists$ ps.-tr. $S$ tel que
\[
  T \prec S \prec T'
\]
et tel que $T$ et $S$, $S$ et $T'$ se raccordent.

De plus, $S$ est unique, et $\partial S = \{a,b\}$. $S$ est formé des $x \in L$ tels que
\[
  T \subset L_{<x}, \quad T' \subset L_{>x} \qquad \text{(cf. b))}
\]

On appelle $S$ le ps.-tr. interstitiel entre $T$ et $T'$.

\textbf{Cor 1} Soit $T$ un tronçon, et supposons que $T$ soit minoré \add{strictement}. Alors l'\uncertain{ens.} des \add{minorants stricts, i.e. des $x \in L$ tels que} $T \subset L_{>x}$ (ens. des minorants stricts) est un \struck{tronçon} ps.-tr. $S$, qui satisfait \struck{\ill{}}
\[
  (*) \qquad S \prec T, \quad S \text{ et } T \text{ se raccordent}, \quad S \text{ sans minorant}
\]
et les conditions \struck{satisf} caractérisent $S$.

\textbf{Cor 2} \struck{Ensemble} On appelle $S$ le \textbf{pstronçon interstitiel gauche} de $T$. On définit de même (si $T$ majoré) le pstr. interstit. droit de $T$.

\page{29}
\struck{\textbf{Définition}}

Soit $\mathrm{Tr}(L)$ l'ens. \struck{ordonné} des ps.-tronçons de $L$, ordonné par $\ll$. Pour toute partie finie tot. ordonnée $\Phi$ de $\mathrm{Tr}(L)$, considérons
\[
  X_{\Phi} = \bigcup_{i \in \Phi} T_i \qquad \text{(NB réunion disjointe)}
\]
\marginal{Prop}
Alors l'application
\[
  \Phi \longmapsto X_{\Phi}, \qquad \underbrace{\mathrm{Drap}^{*}_{\ll}(\mathrm{Tr}(L))}_{\text{drapeaux de } \mathrm{Tr}(L), \text{ y inclus}\ \ldots} \longrightarrow \mathfrak{P}(L)
\]
\note{sous « y inclus », un mot illisible}
est \textbf{injective}.

Les ens. de la forme $\mathcal{U}_{\Phi}$ s'appellent \textbf{modérés}. Pour un tel
\[
  X_{\Phi} = \bigcup T_i
\]
on pose
\[
  \mathrm{Fr}(X_{\Phi}) = \underbrace{\Bigl( \bigcup \partial T_i \Bigr)}_{\partial \mathcal{U}, \text{ drapeau de } L} \smallsetminus \varepsilon
\]
\[
  \varepsilon = \Bigl\{ t \in \partial L \cap \mathcal{U} \;\Big|\; \{t\} \notin \Phi \Bigr\}
\]
i.e. $\{t\}$ n'est pas parmi les $T_i$ ; i.e. $\mathcal{U}$ est un voisinage de $t$ dans $L$ i.e. $\{t\}$ pas \uncertain{isolé} dans $\mathcal{U}$…
\note{dans l'accolade de la page, un mot biffé illisible avant « $\{t\} \notin \Phi$ » ; les « i.e. » qui suivent sont écrits à l'intérieur de l'accolade}
\note{une ligne biffée, « i.e. $\exists x \in L$ … », précède « i.e. $\mathcal{U}$ est un voisinage »}
\marginal{« frontière » relative dans $L$ (mais pas pour la topologie \ill{} de $L$…)}
\marginal{Les $T_i$ s'appellent les « composantes » (connexes !) de \struck{\ill{}} $X$ (mais pas pour la top. canonique) l'ens. de ces composantes \uncertain{on va} noter parfois $\pi_0(X)$}

\page{30}
\struck{Ainsi} Soit $X = X_{\Phi}$ une partie modérée. Un $x \in L$ est dit en posit. rel. modérée avec $X$, s'il l'est avec les $T_i$ \add{(\textbf{composantes} de $X$)} i.e. s'il l'est avec $\bigcup \partial T_i = \partial \mathcal{U}$, ou \uncertain{avec} $\mathrm{Fr}(X)$. Par exemple \struck{si} $x \in X$, ou $x \in \partial \mathcal{U}$, a fortiori $x \in \mathrm{Fr}(X)$, a cette propriété. Soit
\[
  X' = C_{\mathrm{mod}}(\mathcal{U}) = \bigl\{ x \in L \smallsetminus \mathcal{U} \;\big|\; x \text{ est en pos. rel. mod. p.r. à } \mathcal{U} \bigr\}
\]
Alors \struck{$\mathcal{E}$} $X'$ est modérée, et c'est la plus grande partie modérée \struck{$X$} de $L$ telle que \struck{$X$} $Y \subset L \smallsetminus X$, et $Y$ en position relative modérée par rapport à $X$. Si
\[
  \{ T_1 \ll T_2 < \cdots \ll T_n \}
\]
est le drapeau définissant $X$, le drapeau définissant $X'$ est défini par le drapeau formé des $n-1$ tronçons interstitiels \struck{entre} \add{pour} $(T_1, T_2)$, $(T_2, T_3)$, …, $(T_{n-1}, T_n)$, plus le cas échéant (s'ils existent) les tronçons interstitiel gauche de $T_1$, et le tronçon interstitiel droit de $T_n$.

\marginal{Plus généralement, si $X$, $Y$ \struck{\ill{}} modérées, on dit que $X$, $Y$ sont en pos. rel. modérée si \struck{$\partial \mathcal{U}$ et $\partial \mathcal{V}$} $\partial X$, $\partial Y$ le sont, i.e. si \ill{} $\mathrm{Fr}\, X$, $\mathrm{Fr}\, Y$ \ill{} \ill{}. Ex : si $X \subset Y$, alors $(X,Y)$ \ill{} modérée}
\note{note marginale oblique, en partie surchargée}

\page{31}
Si $n'$ est le cardinal du $\Phi' \in \mathrm{Drap}^{\alpha}_{\ll}(\mathrm{Tr}(L))$ qui définit $X'$, on aura
\[
  n' \in \{ n-1, n, n+1 \}
\]
donc
\[
  \bigl| \mathrm{Card}\, \pi_0(X) - \mathrm{Card}\, \pi_0(X') \bigr| \leq 1
\]
De façon précise
\[
  \begin{array}{ll}
    T_1 \text{ non minoré, } T_n \text{ non majoré} & : n' = n-1 \\
    T_1 \text{ minoré et } T_n \text{ non majoré, ou l'inverse} & : n' = n \\
    T_1 \text{ minoré et } T_n \text{ majoré} & : n' = n+1 .
  \end{array}
\]
\struck{(\ill{})}

On considère l'ens. \struck{ordonné} des parties modérées de $L$, $\mathcal{M} = \mathrm{Mod}(L)$. Il y a des structures remarquables

a) Ordre $X \subset Y$

b) Relation : $\{X,Y\}$ modérée \quad (NB $X \subset Y \Rightarrow \{X,Y\}$ modérée)
\marginal{Cette relation est-elle déterminable à l'aide de la relation d'ordre — p. ex. $\{X,Y\}$ mod. ssi $X \wedge Y$ existe ???}

c) Partie $\mathcal{M}_f = \mathrm{Modfer}(L)$ formée des parties modérées dites « fermées » (i.e. $X = X_{\Phi}$, où les $T_i$ ($i \in \Phi$) sont des ps. tr. fermés)

\textbf{Propriétés}

\page{32}
1) Pour la relation d'ordre, il y a un plus petit élément $\emptyset$, et un plus grand élément $L$ (\uncertain{noté} $1$)

2) $X \subset Y \Longrightarrow \{X,Y\}$ modérée

3) Si $\{X,Y\}$ modérée, $X \wedge Y$ et $X \vee Y$ existent (inf et sup)

\struck{3') Si $\{X, Y_i\}$ modérées \ill{} (distributivité) \ill{} mult. mod., $X \wedge (\bigvee Y_i) = \bigvee (X \wedge Y_i)$}

4) \struck{Parmi les} $\forall X \in \mathcal{M}$ \struck{\ill{}}, parmi les $Y \in \mathcal{M}$ tels que $\{X,Y\}$ modérée, et $X \wedge Y = \emptyset$, il y a un plus grand élément $X'$. C'est aussi l'unique élément $X'$ tel que
\[
  \{X', X\} \text{ modérée}, \quad X' \wedge X = \emptyset, \quad X' \vee X = 1
\]
L'application $X \longmapsto X'$ est involutive, renverse la relation d'ordre (antiisomorphisme), \struck{\ill{}}. De plus $\{X', Y'\}$ modérée${}^{*}$ ssi $\{X,Y\}$ mod. et dans \add{ce cas}
\[
  (X \wedge Y)' = X' \vee Y', \quad (X \vee Y)' = X' \wedge Y'
\]

5) $\emptyset, 1 \in \mathcal{M}_f$

6) $\forall X \in \mathcal{M}$, il existe un plus petit \struck{4')} $Y \in \mathcal{M}_f$ tel que $X \subset Y$ (on pose $Y = \overline{X}$).

\struck{(*)} De plus Si $X, Y \in \mathcal{M}$ tels que $\{X,Y\}$ mod., alors $\overline{X \vee Y} = \overline{X} \vee \overline{Y}$

7) Si $\{X,Y\}$ modérée $\{\overline{X}, Y\}$ aussi (donc $\{\overline{X}, \overline{Y}\}$ aussi)

\marginal{Attention $X \vee Y$ \uncertain{n'est} pas l'union $X \cup Y$ !!}
\marginal{Dire \ill{} : Si \struck{\ill{}} $\{X_i\}$ \ill{} famille d'éléments mutuellement modérés, alors ils \uncertain{engendrent} une sous-alg. \ill{} les op. $\wedge$, $\vee$, $'$ \ill{} de Boole, \ill{} finie si $I$ fini}
\marginal{* En fait, $\forall X, Y$ $\{X,Y\}$ mod. ssi $\{X, Y'\}$ mod. (donc ssi $\{X', Y'\}$ mod. \ill{} ssi $\{X', Y\}$ \uncertain{aussi})}
\marginal{Par passage au complément définir $X^{\circ} = (\overline{X'})'$}

\page{33}
8) $\forall X \in \mathcal{M}$, soit $X'$ son « complémentaire ». Comme $\{X, X'\}$ modérée, alors $\{\overline{X}, \overline{X'}\}$ \struck{\ill{}} fermé \ill{}
\[
  \left\{
  \begin{array}{l}
    \mathrm{Fr}\, X \overset{\mathrm{def}}{=} \overline{X} \wedge \overline{X'} \\
    (\mathrm{Fr}\, X)' = X^{\circ} \vee X'^{\circ}
  \end{array}
  \right.
\]
\struck{$\mathrm{Fr}\, X$ est \ill{}} Cela dit $\{X, Y\}$ modérée \struck{\ill{}} $\{\mathrm{Fr}\, X, \mathrm{Fr}\, Y\}$ est modérée $\Longrightarrow$ $\mathrm{Fr}\, X \cup \mathrm{Fr}\, X$ est modérée.
\note{« $\mathrm{Fr}\, X \cup \mathrm{Fr}\, X$ » sur la page}
\marginal{propriété \uncertain{spécifique} au cas des modérés de \uncertain{$L$}, un \ill{} de dimension $1$}

\struck{\textbf{NB} \struck{Si $X'$ est dense} On dit que $X$ est \textbf{rare} si $X^{\circ} = \emptyset$ i.e. $\overline{X'} = 1$ ($X'$ \textbf{dense}). Dans le cas d'espèce, $X$ est rare ssi \ill{} \ill{} \ill{} $X$ est un ens. fini (partie de $L$), ou encore ssi $\mathrm{Fr}\, X$}
\note{ce NB est barré de quatre traits obliques ; il est repris juste après}

Dans le cas d'espèce, conditions équivalentes sur $X \in \mathcal{M}$ :

a) $X \subset \struck{\ill{}}\ \dot{X}$ \quad (i.e. $X^{\circ} = \emptyset$ i.e. $\overline{X'} = 1$ i.e. $X'$ « dense »)

b) $X = \dot{X}$ \quad (i.e. $X$ fermé rare i.e. $X'$ ouvert dense)

c) $X$ est fini en tant que partie de $L$

Conditions sur un ordonné $\mathcal{M}$ satisfaisant les conditions 1) à 7) (avec 4) \ill{} \uncertain{singulièrement} comme il a été dit ; dans la notation algébrique, on a
\[
  \begin{array}{ll}
    X + Y = (X \vee Y) \wedge (X \wedge Y)' & \qquad X \vee Y = X + Y - XY \\
    X \cdot Y = X \wedge Y & \qquad X \wedge Y = X \cdot Y \\
    0 = \text{plus petit él.}, \ 1 = \text{plus grd él.} & \qquad X \leq Y \iff X \cdot Y = X
  \end{array}
\]
\note{la colonne de droite est encadrée, une flèche la rattachant à la ligne « dans la notation algébrique »}

\page{34}
Soit $A \in \mathcal{M}$. On va « relativiser » \struck{\ill{}} par rapport à $A$. Soit \struck{$\mathcal{M}_X = \mathcal{M}_{\leq X}$}
\[
  \mathcal{M}_A = \mathcal{M}_{\leq A} = \{ Y \in \mathcal{M} \mid Y \subset A \}
\]
$\mathcal{M}_A$ muni de la relation d'ordre induite par $\mathcal{M}$, $\{X, Z\}$ est mod. dans $\mathcal{M}_A$ ssi il l'est dans $\mathcal{M}$, $X \in \mathcal{M}_A$ est fermé ssi $X = \overline{X} \wedge A$ i.e. ssi $\exists Z \in \mathcal{M}_f$, $\{Z, A\}$ mod. et $X = Z \wedge A$.

Je dis que les conditions 1) à 7) sont encore satisfaites :

1) $\emptyset_{\mathcal{M}_A} = \emptyset_{\mathcal{M}}$, $1_{\mathcal{M}_A} = X$
\note{« $= X$ » sur la page ; on attend $A$}

2) \struck{\ill{}} trivial

3) item car \struck{\ill{}} $X \overset{\mathcal{M}}{\wedge} Y$ et $X \overset{\mathcal{M}}{\vee} Y$ sont dans $\mathcal{M}_A$ dans ce cas et sont les bornes inf et sup dans $\mathcal{M}_A$.

4) \struck{$X < A$,} $X \in \mathcal{M}_A$, considérons les $Y \in \mathcal{M}_Y$ en position mod. p.r. à $X$, telles que $Y \wedge X = \emptyset$ i.e. \struck{\ill{}} on a $X \subset A$, on considère les $Y \subset A$ tels que $\{X, Y\}$ mod. et $X \wedge Y = \emptyset$. Les $Y$ sont $\leq X'$ (déf. de $X'$) et $\leq A$, donc $\leq X' \wedge A$, qui a les propriétés voulues, et c'est le plus grand de tous, OK.
\note{« $\mathcal{M}_Y$ » sur la page ; on attend $\mathcal{M}_A$}

\page{35}
Considérons des \struck{$X_i$} $X_i \in \mathcal{M}_A$ mutuellement modérés, et les opérations $\complement_A$, $\wedge$, $\vee$ \struck{i.e.} i.e. des $X_i \subset A$ mut. mod. et mod. p.r. à $A$ \uncertain{aussi}) on regarde les op. $\complement_A$ $\wedge$, $\vee$ \struck{(binaires) \ill{}} $\complement_{\mathcal{M}_A}$ (unaire) $\wedge$, $\vee$ (binaires) et on ajoute $0_{\mathcal{M}_A}$, $1_{\mathcal{M}_A}$), on voit que c'est pareil que \struck{\ill{}} de prendre $\complement_{\mathcal{M}}$, $\wedge$, $\vee$, $0_{\mathcal{M}}$, $1_{\mathcal{M}}$ — et de multiplier par $A$. Mais alors on sait que l'on trouve un idéal d'une alg. de Boole, c'est une alg. de Boole…

\struck{7)} Je fais remarquer que $\{X, Y\}$ mod. $\Rightarrow$ $\{\complement_{\mathcal{M}_A} X, \complement_{\mathcal{M}_A} Y\}$ mod., mais $\{\complement_{\mathcal{M}} X, \complement_{\mathcal{M}} Y\}$ \ill{} est mod., donc comme \struck{\ill{}} \ill{} sont mod. avec $A$) $\{A \wedge X', A \wedge Y'\}$ modérée, OK.

5) Trivial

6) Soit $X \subset A$, les fermés dans $A$ sont les $Z = \struck{\ill{}}\ Y \wedge A$, avec $Y \in \mathcal{M}_f$, $\{Y, A\}$ mod. \struck{\ill{}} Le $Z$ qui contient $X$ \struck{\ill{}}, donc $Y$ t.q. $Y \supset X$ i.e. $Y \supset \overline{X}$, et le

\page{36}
plus petit possible est $Y = \overline{X}$, d'où \struck{$Z = \overline{X}$}
\[
  \overline{X}^{(\mathcal{M}_A)} = \overline{X}^{\mathcal{M}} \wedge A
\]
Soient $X, Y \in \mathcal{M}_A$, $\{X, Y\}$ modérée, alors
\[
  \overline{X \vee Y}^{\mathcal{M}_A} = \underbrace{\overline{X \vee Y}^{\mathcal{M}}}_{} \wedge A = (\overline{X} \vee \overline{Y}) \wedge A = \underbrace{(\overline{X} \wedge A)}_{\overline{X}^{A}} \vee (\overline{Y} \wedge A) = \overline{X}^{A} \vee \overline{Y}^{A}
\]
\note{la page dispose ces égalités sur deux lignes, reliées par des accolades}

7) $X, Y \in \mathcal{M}_A$, $\{X, Y\}$ modérée, donc \struck{$\{\overline{X}^{\mathcal{M}}, Y^{\mathcal{M}}\}$ \ill{}} $\{X, Y, A\}$ mod., donc $\{\overline{X}, Y, A\}$ mod., donc $\{\overline{X} \wedge A = \overline{X}^{A}, Y\}$ modérée. \textbf{OK}

\subsection{Composantes}

On dit qu'un $\mathcal{M}$ est \textbf{connexe} si \struck{$\forall X$ fermé et} \struck{ouvert}. $\forall X \in \mathcal{M}$, $X$ fermé et $X$ ouvert $\Longrightarrow X = \emptyset$ ou $X = 1$

\textbf{NB} Soit $X$ tel que $X$ fermé et ouvert, \ill{} qu'il soit mod. \ill{} \ill{}. \ill{} \ill{} ainsi, Alors \struck{LE} $Y$. S'il en est ainsi, les $X$ en question \uncertain{engendrent} une algèbre de Boole. \struck{\ill{}} \struck{\ill{}} \ill{} \ill{} \ill{}

\page{37}
de Boole en question est \textbf{finie}, alors les comp. connexes sont les $X \neq \emptyset$ \textbf{minimaux} qui sont ouverts et fermés. [On suppose $0 \neq 1$ i.e. $\mathcal{M}$ non réduit à un pt — sinon on a \struck{une} aussi une structure « vide »…]

\textbf{Ex} Reprenons $\mathcal{M} = \mathrm{Mod}(L)$, et $\mathcal{M}_A$ où $A \in \mathcal{M} = \mathrm{Mod}(L)$, donc $A$ partie modérée d'un tronçon. On a donc
\[
  A = X_{\Phi} = \bigcup T_i, \qquad \Phi \in \mathrm{Drap}^{*}_{\ll}(\mathrm{Tr}(L))
\]
Je dis que les comp. connexes de $A$ sont les ps.-tronçons composants $T_i$. En effet, on a bien sûr
\[
  A = \bigvee_i T_i \text{ dans } \mathcal{M}, \quad \struck{\text{et}} \text{ et } T_i \wedge T_j = \emptyset \text{ si } i \neq j,
\]
il suffit donc de prouver que a) Chaque $T_i$ est ouvert et fermé dans $A$, et b) Chaque $T_i$ est connexe, i.e. ne peut \uncertain{s'écrire} comme réunion \uncertain{disjointe} de deux parties \struck{\ill{}} disjointes et \add{fermées} $\neq \emptyset$.

\page{38}
Pour a), il suffit de prouver que $T_i$ est fermé dans \struck{\ill{}} $A$ (les $T_j$, $j \neq i$, l'étant aussi, $T_i$ est aussi ouvert…) \struck{\ill{}} i.e. que
\[
  \underbrace{\overline{T_i} \cap A}_{} = T_i ,
\]
\note{sous l'accolade, de sa main : « \ill{} \ill{} modérés bien sûr ! »}
C'est immédiat.

\textbf{NB} Si $T = T_i$ est un ps.-tr., son « adhérence » modérée n'est autre que le ps.-tr. \add{fermé} associé (= $T$ si $T$ réduit à un pt). Il suffit de voir ce que \ill{} \uncertain{quand} $T$ n'est pas fermé, \struck{$a \in \ill{}$ $\mathrm{Fr}(T) = \ill{}$ tel que} un $a \in \partial T \smallsetminus T$ … Si $T = \{a\}$ trivial. Sinon $T$ \ill{} \ill{} \quad \textbf{NB} $(\mathcal{M}_L)_T = \mathcal{M}_T$,

b) $T$ \textbf{connexe} : on regarde $T$ comme un tronçon \uncertain{en lui-même}. \struck{Donc} \ill{} \uncertain{ramène} OPS $T = L$. Mais si $X$ est modéré, \struck{\ill{}} donc
\[
  X = X_{\Phi} = \bigcup_{i \in I} T_i ,
\]
alors \struck{Et} la frontière modérée est la réunion des $\partial T_i$, elle est vide ssi les $\partial T_i$ sont vides. Mais $\partial T_i = \emptyset \Longrightarrow T_i = L$ donc $X = \emptyset$ ($I = \emptyset$) ou $X = L$ ($I \neq \emptyset$) OK.

\page{39}
\textbf{NB} Soit $(X_{\alpha})$ famille \add{finie} de parties modérées de $L$, \struck{mut.} mutuellement modérées, i.e. telles que
\[
  S = \bigcup_{\alpha} \dot{X}_{\alpha}
\]
soit une partie tot. ordonnée : $t_1 < t_2 < \cdots < t_n$. Alors l'alg. de Boole engendrée par les $X_{\alpha}$ \struck{est \ill{}} admet comme \struck{composantes} \struck{parties (\ill{} \ill{} des parties \ill{} comme partie de $L$} (et ps. tr.) « composants » (\ill{} \ill{} \ill{} \ill{}) \ill{} \ill{}
\[
  \underset{\text{si } t_1 \notin \partial L}{[L_{<t_1}]},\ \{t_1\},\ ]t_1, t_2[,\ \{t_2\},\ ]t_1, t_3[,\ \{t_3\}, \ldots, \{t_n\},\ \underset{\text{si } t_n \notin \partial L}{[L_{>t_n}]}
\]
\note{« $]t_1, t_3[$ » tel qu'il se lit ; on attend $]t_2, t_3[$}

Ainsi, cette alg. de Boole est \uncertain{connue} quand on connaît $S$, \uncertain{comme} \ill{} \ill{} $t_1 < \cdots < t_n$ (\ill{} « \uncertain{décomposition} » \add{ou} « subdivision » des tronçons).
\note{les lignes du milieu de la page, avant la liste, sont barrées de deux longs traits et d'une surcharge ; leur lecture est fragmentaire}

\page{40}
\subsection{Caractérisation d'un intervalle ordonné, en tant que réseau}
\note{le titre est de sa main ; « \struck{réseau intervalle} » biffé en tête de la deuxième ligne}

par la donnée de
\[
  S \subset \underset{\displaystyle \mathrm{Drap}_2(L)}{\overset{\|}{\mathfrak{P}_2(L)}}
\]
et pour $\varepsilon = \{a,b\} \in S$, la décomposition de \struck{$L$ en}
\[
  \complement_{\mathrm{mod}}\, \varepsilon = \bigl\{ x \in L \;\big|\; \{a,x\}, \{b,x\} \in S \bigr\}
\]
en « composantes connexes » (par relation $R_{\varepsilon}$)

J'ai commencé à faire l'exercice, mais j'ai l'impression que je n'ai pas la peine de l'expliciter jusqu'au bout. Ça doit avoir été fait depuis longtemps (Hilbert ?) dans le cas totalement ordonné \add{(et divisible)} (en termes de la relation : $s$ est entre $a$ et $b$…)

\end{document}
